\documentclass[11pt]{article}
\usepackage[T1]{fontenc}
\usepackage[utf8]{inputenc}
\usepackage{amsmath,amssymb,amsthm,mathtools}
\usepackage{geometry,graphicx,booktabs,longtable,tabularx,array}
\usepackage{enumitem,needspace,xcolor,tikz,tcolorbox}
\usetikzlibrary{arrows.meta,positioning,calc,fit,backgrounds,decorations.pathreplacing,matrix,patterns}
\tcbuselibrary{breakable}
\usepackage{caption,hyperref,bookmark,cleveref,microtype}
\usepackage{komorebi-publication}
\renewcommand{\KomorebiTitle}{Komorebi Space}
\renewcommand{\KomorebiSubtitle}{Coherent realizations, obstructions, and descent}
\renewcommand{\KomorebiSeries}{Foundations of the Realization Calculus}
\renewcommand{\KomorebiKind}{Monograph}
\renewcommand{\KomorebiEdition}{Monograph edition 1.0.0 / 10 October 2026}
\renewcommand{\KomorebiShortTitle}{Komorebi Space}
\renewcommand{\KomorebiProfile}{monograph}
\KomorebiSetup
\setcounter{tocdepth}{1}
\setcounter{secnumdepth}{2}
\setlist{nosep,leftmargin=1.7em,topsep=4pt}
\hypersetup{pdfkeywords={Komorebi Space,realization calculus,descent,obstruction,Lean,holonomy,closed range}}
\newtheorem{theorem}{Theorem}[section]
\newtheorem{proposition}[theorem]{Proposition}
\newtheorem{lemma}[theorem]{Lemma}
\newtheorem{corollary}[theorem]{Corollary}
\theoremstyle{definition}
\newtheorem{definition}[theorem]{Definition}
\newtheorem{example}[theorem]{Example}
\theoremstyle{remark}
\newtheorem{remark}[theorem]{Remark}
\numberwithin{equation}{section}
\newcommand{\R}{\mathbb R}
\newcommand{\Z}{\mathbb Z}
\newcommand{\N}{\mathbb N}
\newcommand{\K}{\mathbb K}
\newcommand{\V}{\mathcal V}
\newcommand{\W}{\mathcal W}
\newcommand{\id}{\operatorname{id}}
\newcommand{\im}{\operatorname{im}}
\newcommand{\ran}{\operatorname{ran}}
\newcommand{\Hom}{\operatorname{Hom}}
\newcommand{\Ob}{\operatorname{Ob}}
\newcommand{\Ext}{\operatorname{Ext}}
\newcommand{\Match}{\operatorname{Match}}
\newcommand{\Exact}{\operatorname{Exact}}
\newcommand{\dist}{\operatorname{dist}}
\newcommand{\Aut}{\operatorname{Aut}}
\newcommand{\Stab}{\operatorname{Stab}}
\newcommand{\Set}{\mathbf{Set}}
\newcommand{\Cat}{\mathbf{Cat}}
\newcommand{\norm}[1]{\left\lVert#1\right\rVert}
\newcommand{\abs}[1]{\left\lvert#1\right\rvert}
\newcommand{\source}[1]{\par\phantomsection\label{src:#1}{\footnotesize\sffamily\color{KomorebiMuted}Source concordance: \hyperref[app:concordance]{#1}.}\par\smallskip}
\newcommand{\code}[1]{\texttt{\detokenize{#1}}}
\newcommand{\chap}[1]{\par\Needspace{10\baselineskip}\section{#1}}
\newcommand{\partnote}[2]{\clearpage\part{#1}\noindent\begin{minipage}{.92\linewidth}\large #2\end{minipage}\par\vspace{12pt}}
\newtcolorbox{insight}{breakable,colback=KomorebiSoft!55,colframe=KomorebiPine!50,boxrule=.45pt,arc=1pt,left=9pt,right=9pt,top=7pt,bottom=7pt}
\newtcolorbox{scopebox}{breakable,colback=KomorebiWarm,colframe=KomorebiOchre!55,boxrule=.4pt,arc=1pt,left=9pt,right=9pt,top=7pt,bottom=7pt}
\newcommand{\plate}[3]{\begin{figure}[htbp]\centering\input{figures/#1}\caption{#2}\label{#3}\end{figure}}
\tikzset{box/.style={draw=KomorebiPine!70,rounded corners=2pt,fill=KomorebiSoft!50,align=center,inner sep=8pt,font=\small},arr/.style={-{Latex[length=2.2mm]},draw=KomorebiPine,thick},thinarr/.style={-{Latex[length=1.8mm]},draw=KomorebiMuted},lab/.style={font=\small,align=center},dot/.style={circle,fill=KomorebiPine,inner sep=2.3pt}}
\renewcommand{\arraystretch}{1.18}
\raggedbottom
\widowpenalty=7000\clubpenalty=7000
\setlength{\emergencystretch}{3em}

% Compact only the contents entries, without reducing body type size.
\usepackage{etoolbox}
\makeatletter
\patchcmd{\l@section}{1.0em}{0.6em}{}{}
\makeatother

\renewcommand{\KomorebiKind}{Lean and evidence guide}
\renewcommand{\KomorebiSubtitle}{Sources, theorem correspondence, and verification}
\renewcommand{\KomorebiEdition}{Companion edition 1.0.0 / 10 October 2026}
\begin{document}
\KomorebiCover
\section*{What this guide establishes}
This is a source and evidence guide to the monograph \emph{Komorebi Space: Coherent realizations, obstructions, and descent}. It connects the printed mathematics to actual recovered Lean declarations, preserves the lineage of the checked sources, and states which checks were run for this publication.

The publication contains 54 principal exposition items with stable identifiers KS01 through KS54. Fifty-three items have source correspondences, comprising 112 declaration references. These are references, not a count of distinct new theorems. Several printed statements assemble related declarations, and some examples add an explanatory derivation. KS54 is a separately identified written example with a printed proof, not a new Lean theorem.

\begin{scopebox}
\textbf{Current execution boundary}\par
The source and archive integrity checks, correspondence resolver, independent exact examples, recovered finite checker, and PDF build were run for this edition. A fresh Lean compiler run was \textbf{not} performed. No Lean executable was available in the working environment, and the external pinned runtime was not restored. Historical compiler reports are retained as historical evidence, not relabeled as a current build.
\end{scopebox}

\subsection*{Four questions that should not be collapsed}
A source declaration can be present without a fresh build. A successful build can check a theorem with narrower hypotheses than an edited paragraph suggests. A theorem can be correct without its application hypotheses being established. A correct theorem can also be classical rather than historically new.

For that reason, this guide separates source correspondence, historical compilation, fresh publication checks, and the mathematical scope of the monograph. The complete source types remain authoritative for formal scope. The printed proofs and the correspondence notes make the editorial interpretation inspectable.

\subsection*{Files to start with}
The machine-readable source map is \path{data/claim_ledger.json}. The shorter spreadsheet-compatible form is \path{data/claim_ledger.csv}. The original archives are in \path{evidence/archives/}. Current execution results are in \path{audit/}. The root progress ledger records completed work, the known limitations, and the next reproducible restart point.

\clearpage\section{The recovered source lineages}
\subsection{F: the cumulative foundation archive}
\path{Komorebi_Space_Lean_v0.20.0.zip} is a complete recovered source release. It contains 122 production proof modules and a root module. Its historical inventory reports 1,299 named theorem declarations and 696 indexed constructions. These counts include auxiliary lemmas and definitions.

This tree contains coherent views, partial extension, additive obstruction, tower repair, residual cascades, nonabelian cocycles, effective groupoid descent, the composition-site stack theorem, equivalent presentation results, and finite-dimensional amplified holonomy repair.

The historical evidence distinguishes a full cumulative source build from the subsequent archive replay. The replay reuses an already checked 117-module baseline and recompiles the five new modules and root. It compares all 123 source and compiled-module hashes. It is not a second clean build of every foundation module.

\subsection{M: the focused moving-space archive}
\path{Komorebi_Moving_Exact_Spaces_v0.32.0.zip} is a focused source closure containing 11 production modules and a root. Its inventory reports 151 named theorem declarations and 47 constructions. Six modules are inherited dependencies of the focus, and five are additions.

Its recorded checks include two clean builds of this supplied closure, declaration axiom requests, and negative controls. The five false proof files must fail. The sixth control introduces a custom axiom deliberately: it compiles, but the subsequent axiom audit must detect the unapproved assumption. That control is not described as a compiler rejection.

This focused source tree proves closed-range repair criteria, moving-kernel bounds, derived perturbed splittings, explicit kernel transport, and the bridge to the existing quantitative-descent repair. Its higher version number does not make it a replacement for the whole foundation archive.

\subsection{S: a separately recovered source}
\path{evidence/selected/HilbertHolonomyRepair.lean} implements the Hilbert-fiber amplified holonomy repair and its closed-range necessity criterion. The file is retained with its hash and exact declaration correspondence. This edition does not certify its complete build closure or apply the focused archive's replay status to this separate source.

\begin{insight}
F, M, and S are labels for provenance. A theorem combining their mathematical ideas is not thereby a newly compiled union of their source trees. The monograph's finite and infinite-dimensional holonomy discussion identifies the difference explicitly.
\end{insight}

\clearpage\section{The pinned environment and the trust boundary}
The archived records identify Lean 4.34.0-rc2 and mathlib revision
\begin{center}\small\path{aeb72f1e9981b6d35b426b16dca7759be3e7ba06}.\end{center}
The permitted standard axioms in the historical reports are \code{propext}, \code{Classical.choice}, and \code{Quot.sound}. These are not project-specific axioms asserting the desired mathematical conclusion.

The kernel, compiler binary, and imported cached dependencies remain within the historical trust boundary. A clean project build is not an independent rebuild of Lean or of every mathlib dependency. Archive hashes identify the delivered bytes. They are not external attestations of the person or machine that originally executed the compiler.

\subsection{What was checked again for this edition}
The publication integrity script recomputed the hashes and CRC checks of four frozen archives. It matched every extracted production source and root against the frozen source archive. It also matched all 123 foundation source hashes and all 12 focused source hashes against their retained historical replay records.

The local KomorebiSpace import closure is present inside each complete source tree. A lexical scan, after removing comments and strings, found no \code{sorry}, \code{admit}, \code{native_decide}, project-specific \code{axiom} declaration, or \code{unsafe} declaration in those production sources. This scan does not replace an imported axiom-dependency audit.

The correspondence resolver checked the 112 named references against the frozen source inventories and actual files. It recorded line numbers and source hashes. It does not mechanically prove that every word of the edited prose is equivalent to a Lean type.

\subsection{Fresh exact arithmetic}
The independent publication checker executed \textbf{14,224 assertions}. Its deterministic families include 201 integer seeds, nonabelian \(S_3\) gauge examples, three-view and four-view covers, 43 rational moving-line slopes, 243 short residual vectors, five finite cascade sizes, 20 amplified chains, and nested projection examples.

The recovered focused finite checker was also executed again. It passed \textbf{2,767 assertions} over 36 split systems, 324 candidate vectors and measurement cases, 17 moving-line slopes, and 32 rank-loss witnesses. Its own manifest validator passed. These two checks use exact integers, permutations, and rational arithmetic rather than floating-point tolerance as their source of equality.

Together these are \textbf{16,991 executed finite assertions}, not 16,991 theorems. They check implementations and examples. They do not establish arbitrary-category descent, arbitrary-dimensional operator theorems, or a new compiler run.

\clearpage\section{Reading three important formal interfaces}
\subsection{The standard stack target}
The foundation declaration
\begin{center}\small\path{KomorebiSpace.GroupViewSystem.patch_isStack}\end{center}
is located in \path{CompositionStack.lean}. Its target is the actual mathlib predicate \code{Pseudofunctor.IsStack}, applied to the constructed patch pseudofunctor and composition topology. It is not a custom predicate defined to repeat a desired conclusion.

The proof uses a constructed equivalence of descent categories. Its essential-surjectivity part supplies a global cocycle with the specified local comparison gauges. Its full-faithfulness part glues and uniquely recovers compatible gauge arrows. The composition-cover condition provides a patch in which each multiplication equation can be checked.

The source does not establish stackhood for arbitrary coefficients or an arbitrary site. It also does not identify every semantic Komorebi model with this cocycle stack.

\subsection{The closed-range criterion}
The focused declaration
\begin{center}\small\path{KomorebiSpace.HilbertKernelProjection.closed_range_iff_bound}\end{center}
quantifies over the actual bounded linear defect map. In the complete-target setting, closed range is equivalent to a finite uniform bound of projection error by residual. The orthogonal projection itself is constructed before closed range is assumed.

The companion declaration \path{closed_range_iff_near_exact} allows an arbitrary choice of exact comparison point for each input. This matters: a nonlinear selection cannot bypass the failure of the same uniform estimate. A claim that exact repair always exists must not be read as a claim that it always has controlled residual gain.

\subsection{The moving-kernel estimate}
The focused declaration
\begin{center}\small\path{KomorebiSpace.MovingKernel.projection_opNorm}\end{center}
compares the two original kernel projectors. Its hypotheses include one residual certificate for each defect. It does not assume that the kernels are equal or that one includes the other. The bound is
\[
 \norm{q-q'}\le\max(C,C')\norm{D-D'}.
\]
A separate split-surjective theorem constructs the new certificate inside a strict perturbation radius. The general two-sided theorem should not be substituted for that construction when only one certificate is known.

\clearpage\section{The correspondence catalogue}
Each entry below records the editorial role, the actual source declaration, its starting line in the delivered source, and a short source-hash prefix. The JSON ledger contains the full hash and source path. Repeated references are intentional when one implementation supports more than one part of the exposition.

\textbf{F} is the foundation tree. \textbf{M} is the focused moving-space tree. \textbf{S} is the separate recovered source. \textbf{W} is a written addition. A catalogue entry does not upgrade an illustrative derivation into a newly compiled theorem.
\par\noindent\begin{minipage}{\linewidth}\subsection*{KS01 / Limit universal property}
Grouped source statements. The universe and smallness conventions are retained.
\par\smallskip\noindent{\small\textbf{F} \path{ViewSections.lean}, line 39.\par\noindent\path{KomorebiSpace.ViewSystem.evaluations_jointly_injective}\par{\sffamily\footnotesize Source SHA-256 prefix: 53bb5c493e5cada2.}}\par
\par\smallskip\noindent{\small\textbf{F} \path{ViewSections.lean}, line 99.\par\noindent\path{KomorebiSpace.ViewSystem.Cone.exists_unique_lift}\par{\sffamily\footnotesize Source SHA-256 prefix: 53bb5c493e5cada2.}}\par
\smallskip\noindent{\color{KomorebiMuted}\rule{\linewidth}{.3pt}}\par\end{minipage}\par\medskip
\par\noindent\begin{minipage}{\linewidth}\subsection*{KS02 / Functorial transport}
Direct coordinate-level transport laws.
\par\smallskip\noindent{\small\textbf{F} \path{ViewSections.lean}, line 73.\par\noindent\path{KomorebiSpace.ViewSystem.Morphism.onAmbient_id}\par{\sffamily\footnotesize Source SHA-256 prefix: 53bb5c493e5cada2.}}\par
\par\smallskip\noindent{\small\textbf{F} \path{ViewSections.lean}, line 75.\par\noindent\path{KomorebiSpace.ViewSystem.Morphism.onAmbient_comp}\par{\sffamily\footnotesize Source SHA-256 prefix: 53bb5c493e5cada2.}}\par
\smallskip\noindent{\color{KomorebiMuted}\rule{\linewidth}{.3pt}}\par\end{minipage}\par\medskip
\par\noindent\begin{minipage}{\linewidth}\subsection*{KS03 / Terminal evaluation}
An already-global terminal fiber is recovered, not proved inhabited.
\par\smallskip\noindent{\small\textbf{F} \path{ViewSections.lean}, line 142.\par\noindent\path{KomorebiSpace.ViewSystem.terminalEvaluationEquiv}\par{\sffamily\footnotesize Source SHA-256 prefix: 53bb5c493e5cada2.}}\par
\smallskip\noindent{\color{KomorebiMuted}\rule{\linewidth}{.3pt}}\par\end{minipage}\par\medskip
\par\noindent\begin{minipage}{\linewidth}\subsection*{KS04 / Strict section gluing}
Requires a transition cover, not just object coverage.
\par\smallskip\noindent{\small\textbf{F} \path{SectionGluing.lean}, line 98.\par\noindent\path{KomorebiSpace.ViewSystem.gluingEquiv}\par{\sffamily\footnotesize Source SHA-256 prefix: cf412d554b19ca22.}}\par
\par\smallskip\noindent{\small\textbf{F} \path{SectionGluing.lean}, line 85.\par\noindent\path{KomorebiSpace.ViewSystem.Matching.exists_unique_glue}\par{\sffamily\footnotesize Source SHA-256 prefix: cf412d554b19ca22.}}\par
\smallskip\noindent{\color{KomorebiMuted}\rule{\linewidth}{.3pt}}\par\end{minipage}\par\medskip
\par\noindent\begin{minipage}{\linewidth}\subsection*{KS05 / Object-only cover failure}
The Boolean counterexample is implemented.
\par\smallskip\noindent{\small\textbf{F} \path{ViewBoundaryExamples.lean}, line 36.\par\noindent\path{KomorebiSpace.ViewBoundaryExamples.singleton_matching_not_gluable}\par{\sffamily\footnotesize Source SHA-256 prefix: c094eb03ea6b155f.}}\par
\smallskip\noindent{\color{KomorebiMuted}\rule{\linewidth}{.3pt}}\par\end{minipage}\par\medskip
\par\noindent\begin{minipage}{\linewidth}\subsection*{KS06 / Correspondence fibers}
Ordinary set fibers, not homotopy fibers.
\par\smallskip\noindent{\small\textbf{F} \path{ViewCorrespondence.lean}, line 43.\par\noindent\path{KomorebiSpace.ViewSystem.bridgeEquiv}\par{\sffamily\footnotesize Source SHA-256 prefix: a89331eabd7662f8.}}\par
\par\smallskip\noindent{\small\textbf{F} \path{ViewCorrespondence.lean}, line 94.\par\noindent\path{KomorebiSpace.ViewSystem.bridge_exists_iff_matching}\par{\sffamily\footnotesize Source SHA-256 prefix: a89331eabd7662f8.}}\par
\smallskip\noindent{\color{KomorebiMuted}\rule{\linewidth}{.3pt}}\par\end{minipage}\par\medskip
\par\noindent\begin{minipage}{\linewidth}\subsection*{KS07 / Nonempty fibers without a section}
The one-object Boolean swap has no global fixed choice.
\par\smallskip\noindent{\small\textbf{F} \path{ViewBoundaryExamples.lean}, line 90.\par\noindent\path{KomorebiSpace.ViewBoundaryExamples.swap_no_ambient}\par{\sffamily\footnotesize Source SHA-256 prefix: c094eb03ea6b155f.}}\par
\par\smallskip\noindent{\small\textbf{F} \path{ViewBoundaryExamples.lean}, line 87.\par\noindent\path{KomorebiSpace.ViewBoundaryExamples.swap_transports_bijective}\par{\sffamily\footnotesize Source SHA-256 prefix: c094eb03ea6b155f.}}\par
\smallskip\noindent{\color{KomorebiMuted}\rule{\linewidth}{.3pt}}\par\end{minipage}\par\medskip
\par\noindent\begin{minipage}{\linewidth}\subsection*{KS08 / Additive view differential}
Uses the actual arrow-indexed differential.
\par\smallskip\noindent{\small\textbf{F} \path{SectionObstruction.lean}, line 58.\par\noindent\path{KomorebiSpace.AdditiveViewSystem.differential_comp}\par{\sffamily\footnotesize Source SHA-256 prefix: 9089dd26b6bf34a0.}}\par
\par\smallskip\noindent{\small\textbf{F} \path{SectionObstruction.lean}, line 64.\par\noindent\path{KomorebiSpace.AdditiveViewSystem.differential_zero_iff}\par{\sffamily\footnotesize Source SHA-256 prefix: 9089dd26b6bf34a0.}}\par
\smallskip\noindent{\color{KomorebiMuted}\rule{\linewidth}{.3pt}}\par\end{minipage}\par\medskip
\par\noindent\begin{minipage}{\linewidth}\subsection*{KS09 / Additive extension criterion}
The quotient is defined independently of solvability.
\par\smallskip\noindent{\small\textbf{F} \path{SectionObstruction.lean}, line 126.\par\noindent\path{KomorebiSpace.AdditiveViewSystem.extensionClass_zero_iff}\par{\sffamily\footnotesize Source SHA-256 prefix: 9089dd26b6bf34a0.}}\par
\par\smallskip\noindent{\small\textbf{F} \path{SectionObstruction.lean}, line 142.\par\noindent\path{KomorebiSpace.AdditiveViewSystem.extensionClass_independent}\par{\sffamily\footnotesize Source SHA-256 prefix: 9089dd26b6bf34a0.}}\par
\smallskip\noindent{\color{KomorebiMuted}\rule{\linewidth}{.3pt}}\par\end{minipage}\par\medskip
\par\noindent\begin{minipage}{\linewidth}\subsection*{KS10 / Extension ambiguity torsor}
Translation requires one chosen extension.
\par\smallskip\noindent{\small\textbf{F} \path{SectionExtensionFibers.lean}, line 35.\par\noindent\path{KomorebiSpace.AdditiveViewSystem.extensionParameters}\par{\sffamily\footnotesize Source SHA-256 prefix: e422b1f42cdb28a1.}}\par
\par\smallskip\noindent{\small\textbf{F} \path{SectionExtensionFibers.lean}, line 62.\par\noindent\path{KomorebiSpace.AdditiveViewSystem.unique_extension_iff}\par{\sffamily\footnotesize Source SHA-256 prefix: e422b1f42cdb28a1.}}\par
\smallskip\noindent{\color{KomorebiMuted}\rule{\linewidth}{.3pt}}\par\end{minipage}\par\medskip
\par\noindent\begin{minipage}{\linewidth}\subsection*{KS11 / Parity seed example}
The disconnected free-coordinate variant in the text is an additional written consequence, not a separate Lean declaration.
\par\smallskip\noindent{\small\textbf{F} \path{SectionExtensionExample.lean}, line 65.\par\noindent\path{KomorebiSpace.SectionExtensionExample.extensionClass_zero_iff_even}\par{\sffamily\footnotesize Source SHA-256 prefix: 12f3be43a2401a41.}}\par
\par\smallskip\noindent{\small\textbf{F} \path{SectionExtensionExample.lean}, line 79.\par\noindent\path{KomorebiSpace.SectionExtensionExample.odd_seed_obstructed}\par{\sffamily\footnotesize Source SHA-256 prefix: 12f3be43a2401a41.}}\par
\par\smallskip\noindent{\small\textbf{F} \path{SectionExtensionExample.lean}, line 85.\par\noindent\path{KomorebiSpace.SectionExtensionExample.nonzero_seed_extends}\par{\sffamily\footnotesize Source SHA-256 prefix: 12f3be43a2401a41.}}\par
\smallskip\noindent{\color{KomorebiMuted}\rule{\linewidth}{.3pt}}\par\end{minipage}\par\medskip
\par\noindent\begin{minipage}{\linewidth}\subsection*{KS12 / Exactness at partial sections}
No surjectivity onto the whole obstruction quotient is asserted.
\par\smallskip\noindent{\small\textbf{F} \path{SectionExactSequence.lean}, line 108.\par\noindent\path{KomorebiSpace.AdditiveViewSystem.connecting_kernel_eq_range}\par{\sffamily\footnotesize Source SHA-256 prefix: 242e72f1d4e431ee.}}\par
\par\smallskip\noindent{\small\textbf{F} \path{SectionExactSequence.lean}, line 132.\par\noindent\path{KomorebiSpace.AdditiveViewSystem.restriction_surjective_iff_connecting_zero}\par{\sffamily\footnotesize Source SHA-256 prefix: 242e72f1d4e431ee.}}\par
\smallskip\noindent{\color{KomorebiMuted}\rule{\linewidth}{.3pt}}\par\end{minipage}\par\medskip
\par\noindent\begin{minipage}{\linewidth}\subsection*{KS13 / Extension by stages}
The dependent sum retains the intermediate choice.
\par\smallskip\noindent{\small\textbf{F} \path{SectionExtensionStages.lean}, line 28.\par\noindent\path{KomorebiSpace.ViewSystem.extensionStages}\par{\sffamily\footnotesize Source SHA-256 prefix: 8196e1272a7300f7.}}\par
\par\smallskip\noindent{\small\textbf{F} \path{SectionExtensionStages.lean}, line 50.\par\noindent\path{KomorebiSpace.ViewSystem.extension_nonempty_iff_intermediate}\par{\sffamily\footnotesize Source SHA-256 prefix: 8196e1272a7300f7.}}\par
\smallskip\noindent{\color{KomorebiMuted}\rule{\linewidth}{.3pt}}\par\end{minipage}\par\medskip
\par\noindent\begin{minipage}{\linewidth}\subsection*{KS14 / Protected correction}
The protected observable is a declared additive map.
\par\smallskip\noindent{\small\textbf{F} \path{Correction.lean}, line 42.\par\noindent\path{KomorebiSpace.CorrectionProblem.flatten_iff}\par{\sffamily\footnotesize Source SHA-256 prefix: 089d9155e0a43f66.}}\par
\par\smallskip\noindent{\small\textbf{F} \path{Correction.lean}, line 56.\par\noindent\path{KomorebiSpace.CorrectionProblem.obstruction_invariant}\par{\sffamily\footnotesize Source SHA-256 prefix: 089d9155e0a43f66.}}\par
\smallskip\noindent{\color{KomorebiMuted}\rule{\linewidth}{.3pt}}\par\end{minipage}\par\medskip
\par\noindent\begin{minipage}{\linewidth}\subsection*{KS15 / Obstruction witness}
A witness annihilates all admissible correction defects.
\par\smallskip\noindent{\small\textbf{F} \path{Correction.lean}, line 64.\par\noindent\path{KomorebiSpace.CorrectionProblem.nonzero_certificate}\par{\sffamily\footnotesize Source SHA-256 prefix: 089d9155e0a43f66.}}\par
\smallskip\noindent{\color{KomorebiMuted}\rule{\linewidth}{.3pt}}\par\end{minipage}\par\medskip
\par\noindent\begin{minipage}{\linewidth}\subsection*{KS16 / Exact tower maps and realized points}
The maps use a fixed sequential schedule. The record instances also supply identity and composition.
\par\smallskip\noindent{\small\textbf{F} \path{Towers.lean}, line 54.\par\noindent\path{KomorebiSpace.Tower.PreHom.exact_comp}\par{\sffamily\footnotesize Source SHA-256 prefix: fdd2819fff6150d2.}}\par
\par\smallskip\noindent{\small\textbf{F} \path{Realization.lean}, line 66.\par\noindent\path{KomorebiSpace.Resolution.Map.onLimit_comp}\par{\sffamily\footnotesize Source SHA-256 prefix: 90e0a22f2f7eebbe.}}\par
\smallskip\noindent{\color{KomorebiMuted}\rule{\linewidth}{.3pt}}\par\end{minipage}\par\medskip
\par\noindent\begin{minipage}{\linewidth}\subsection*{KS17 / Tower differential composition}
No additional mixed term is present in this differential.
\par\smallskip\noindent{\small\textbf{F} \path{Towers.lean}, line 75.\par\noindent\path{KomorebiSpace.Tower.PreHom.defect_comp}\par{\sffamily\footnotesize Source SHA-256 prefix: fdd2819fff6150d2.}}\par
\par\smallskip\noindent{\small\textbf{F} \path{TowerCorrection.lean}, line 39.\par\noindent\path{KomorebiSpace.TowerCorrection.exact_iff_differential_zero}\par{\sffamily\footnotesize Source SHA-256 prefix: 6710b485edbbb676.}}\par
\smallskip\noindent{\color{KomorebiMuted}\rule{\linewidth}{.3pt}}\par\end{minipage}\par\medskip
\par\noindent\begin{minipage}{\linewidth}\subsection*{KS18 / Cofinal point equivalence}
The subsequent alternating example is an explanatory specialization, not a statement that residual norms are invariant.
\par\smallskip\noindent{\small\textbf{F} \path{Cofinal.lean}, line 94.\par\noindent\path{KomorebiSpace.Resolution.cofinalEquiv}\par{\sffamily\footnotesize Source SHA-256 prefix: 0306f5a0a7f05e14.}}\par
\smallskip\noindent{\color{KomorebiMuted}\rule{\linewidth}{.3pt}}\par\end{minipage}\par\medskip
\par\noindent\begin{minipage}{\linewidth}\subsection*{KS19 / Geometric operator repair}
The correction map and residual contraction are hypotheses.
\par\smallskip\noindent{\small\textbf{F} \path{RepairDynamics.lean}, line 70.\par\noindent\path{KomorebiSpace.RepairDynamics.exactification_from_correction_operator}\par{\sffamily\footnotesize Source SHA-256 prefix: 419b196d29b2f14d.}}\par
\par\smallskip\noindent{\small\textbf{F} \path{RepairProjection.lean}, line 128.\par\noindent\path{KomorebiSpace.RepairDynamics.Scheme.projection_bound}\par{\sffamily\footnotesize Source SHA-256 prefix: 4292822995be802a.}}\par
\par\smallskip\noindent{\small\textbf{F} \path{RepairProjection.lean}, line 132.\par\noindent\path{KomorebiSpace.RepairDynamics.Scheme.projection_idempotent}\par{\sffamily\footnotesize Source SHA-256 prefix: 4292822995be802a.}}\par
\smallskip\noindent{\color{KomorebiMuted}\rule{\linewidth}{.3pt}}\par\end{minipage}\par\medskip
\par\noindent\begin{minipage}{\linewidth}\subsection*{KS20 / Constructed split tower correction}
Target splittings and the strict propagation gap are retained.
\par\smallskip\noindent{\small\textbf{F} \path{SplitTowerRepair.lean}, line 154.\par\noindent\path{KomorebiSpace.BanachTower.Splitting.differential_correction}\par{\sffamily\footnotesize Source SHA-256 prefix: 2dca3052b0338605.}}\par
\par\smallskip\noindent{\small\textbf{F} \path{SplitTowerRepair.lean}, line 199.\par\noindent\path{KomorebiSpace.BanachTower.Splitting.repair_error}\par{\sffamily\footnotesize Source SHA-256 prefix: 2dca3052b0338605.}}\par
\smallskip\noindent{\color{KomorebiMuted}\rule{\linewidth}{.3pt}}\par\end{minipage}\par\medskip
\par\noindent\begin{minipage}{\linewidth}\subsection*{KS21 / Split semantic and residual coordinates}
A continuous linear equivalence, not an isometry or an absolute semantic coordinate system.
\par\smallskip\noindent{\small\textbf{F} \path{SplitTowerCoordinates.lean}, line 98.\par\noindent\path{KomorebiSpace.BanachTower.Splitting.coordinates}\par{\sffamily\footnotesize Source SHA-256 prefix: c028e9936850deb9.}}\par
\par\smallskip\noindent{\small\textbf{F} \path{SplitTowerCoordinates.lean}, line 68.\par\noindent\path{KomorebiSpace.BanachTower.Splitting.coordinate_inverse}\par{\sffamily\footnotesize Source SHA-256 prefix: c028e9936850deb9.}}\par
\par\smallskip\noindent{\small\textbf{F} \path{SplitTowerSemantics.lean}, line 141.\par\noindent\path{KomorebiSpace.BanachTower.Splitting.residual_observable_injective}\par{\sffamily\footnotesize Source SHA-256 prefix: ecb64d5ac3d5369d.}}\par
\smallskip\noindent{\color{KomorebiMuted}\rule{\linewidth}{.3pt}}\par\end{minipage}\par\medskip
\par\noindent\begin{minipage}{\linewidth}\subsection*{KS22 / Growing identity repair}
The source bonding maps have the factor one half.
\par\smallskip\noindent{\small\textbf{F} \path{GrowingTower.lean}, line 195.\par\noindent\path{KomorebiSpace.GrowingTower.identityRepair_formula}\par{\sffamily\footnotesize Source SHA-256 prefix: 988f6ae0da7405cd.}}\par
\smallskip\noindent{\color{KomorebiMuted}\rule{\linewidth}{.3pt}}\par\end{minipage}\par\medskip
\par\noindent\begin{minipage}{\linewidth}\subsection*{KS23 / A bounded residual without bounded primitive}
This is a counterexample at the threshold, not necessity of the gap in every tower.
\par\smallskip\noindent{\small\textbf{F} \path{IdentityTowerObstruction.lean}, line 55.\par\noindent\path{KomorebiSpace.IdentityTowerObstruction.no_bounded_primitive}\par{\sffamily\footnotesize Source SHA-256 prefix: 026f9504c3c35507.}}\par
\par\smallskip\noindent{\small\textbf{F} \path{IdentityTowerObstruction.lean}, line 76.\par\noindent\path{KomorebiSpace.IdentityTowerObstruction.repair_exists_but_right_inverse_does_not}\par{\sffamily\footnotesize Source SHA-256 prefix: 026f9504c3c35507.}}\par
\smallskip\noindent{\color{KomorebiMuted}\rule{\linewidth}{.3pt}}\par\end{minipage}\par\medskip
\par\noindent\begin{minipage}{\linewidth}\subsection*{KS24 / Dependence on target splitting}
The two first-level matrices in the exposition are derived from the displayed source recurrence.
\par\smallskip\noindent{\small\textbf{F} \path{SplittingDependence.lean}, line 101.\par\noindent\path{KomorebiSpace.SplittingDependence.repair_operators_differ}\par{\sffamily\footnotesize Source SHA-256 prefix: b5aa369adbd67f9f.}}\par
\par\smallskip\noindent{\small\textbf{F} \path{SplittingDependence.lean}, line 112.\par\noindent\path{KomorebiSpace.SplittingDependence.protected_profiles_differ}\par{\sffamily\footnotesize Source SHA-256 prefix: b5aa369adbd67f9f.}}\par
\smallskip\noindent{\color{KomorebiMuted}\rule{\linewidth}{.3pt}}\par\end{minipage}\par\medskip
\par\noindent\begin{minipage}{\linewidth}\subsection*{KS25 / Lossless dyadic residual split}
Maps are the actual typed dyadic coarsener and its supplied-section lift.
\par\smallskip\noindent{\small\textbf{F} \path{DyadicResidualSplit.lean}, line 39.\par\noindent\path{KomorebiSpace.DyadicResidual.coarsen_pair}\par{\sffamily\footnotesize Source SHA-256 prefix: c4d40d1c7797e5ad.}}\par
\par\smallskip\noindent{\small\textbf{F} \path{DyadicResidualSplit.lean}, line 113.\par\noindent\path{KomorebiSpace.DyadicResidual.coarsen_lift}\par{\sffamily\footnotesize Source SHA-256 prefix: c4d40d1c7797e5ad.}}\par
\par\smallskip\noindent{\small\textbf{F} \path{DyadicResidualSplit.lean}, line 153.\par\noindent\path{KomorebiSpace.DyadicResidual.reconstruction}\par{\sffamily\footnotesize Source SHA-256 prefix: c4d40d1c7797e5ad.}}\par
\smallskip\noindent{\color{KomorebiMuted}\rule{\linewidth}{.3pt}}\par\end{minipage}\par\medskip
\par\noindent\begin{minipage}{\linewidth}\subsection*{KS26 / Cascade telescope and persistent part}
Generic split cascade. No norm convergence is inferred.
\par\smallskip\noindent{\small\textbf{F} \path{CascadeProjections.lean}, line 79.\par\noindent\path{KomorebiSpace.CoarseningCascade.remainder_nested_left}\par{\sffamily\footnotesize Source SHA-256 prefix: b2865b41336df871.}}\par
\par\smallskip\noindent{\small\textbf{F} \path{CascadeProjections.lean}, line 84.\par\noindent\path{KomorebiSpace.CoarseningCascade.remainder_nested_right}\par{\sffamily\footnotesize Source SHA-256 prefix: b2865b41336df871.}}\par
\par\smallskip\noindent{\small\textbf{F} \path{CascadeProjections.lean}, line 136.\par\noindent\path{KomorebiSpace.CoarseningCascade.persistent_iff_details_zero}\par{\sffamily\footnotesize Source SHA-256 prefix: b2865b41336df871.}}\par
\par\smallskip\noindent{\small\textbf{F} \path{CoarseningCascade.lean}, line 73.\par\noindent\path{KomorebiSpace.CoarseningCascade.reconstruction}\par{\sffamily\footnotesize Source SHA-256 prefix: 2719d5a3accb4dde.}}\par
\smallskip\noindent{\color{KomorebiMuted}\rule{\linewidth}{.3pt}}\par\end{minipage}\par\medskip
\par\noindent\begin{minipage}{\linewidth}\subsection*{KS27 / Recovery by density}
Uniform boundedness of the restored remainders is a hypothesis.
\par\smallskip\noindent{\small\textbf{F} \path{CascadeDensity.lean}, line 75.\par\noindent\path{KomorebiSpace.CoarseningCascade.recoverable_iff_closure_finiteCore}\par{\sffamily\footnotesize Source SHA-256 prefix: 11d90a5995354d99.}}\par
\par\smallskip\noindent{\small\textbf{F} \path{CascadeDensity.lean}, line 81.\par\noindent\path{KomorebiSpace.CoarseningCascade.all_recoverable_iff_dense}\par{\sffamily\footnotesize Source SHA-256 prefix: 11d90a5995354d99.}}\par
\smallskip\noindent{\color{KomorebiMuted}\rule{\linewidth}{.3pt}}\par\end{minipage}\par\medskip
\par\noindent\begin{minipage}{\linewidth}\subsection*{KS28 / Dyadic finite-coordinate recovery}
The printed proof is explicitly an outline. The source contains the dependent-index induction.
\par\smallskip\noindent{\small\textbf{F} \path{DyadicInfiniteReconstruction.lean}, line 71.\par\noindent\path{KomorebiSpace.DyadicResidual.System.infinite_details_injective}\par{\sffamily\footnotesize Source SHA-256 prefix: 0c81242c3c6c0715.}}\par
\par\smallskip\noindent{\small\textbf{F} \path{DyadicInfiniteReconstruction.lean}, line 201.\par\noindent\path{KomorebiSpace.DyadicResidual.System.reconstruction_coordinate}\par{\sffamily\footnotesize Source SHA-256 prefix: 0c81242c3c6c0715.}}\par
\smallskip\noindent{\color{KomorebiMuted}\rule{\linewidth}{.3pt}}\par\end{minipage}\par\medskip
\par\noindent\begin{minipage}{\linewidth}\subsection*{KS29 / Shift cascade and c0}
The scalar expanding-lift boundary is separately implemented in CascadeBoundaryExamples.
\par\smallskip\noindent{\small\textbf{F} \path{SequenceCascade.lean}, line 140.\par\noindent\path{KomorebiSpace.SequenceCascade.recoverable_iff_vanishes}\par{\sffamily\footnotesize Source SHA-256 prefix: 0a384fa460f07336.}}\par
\par\smallskip\noindent{\small\textbf{F} \path{SequenceCascade.lean}, line 194.\par\noindent\path{KomorebiSpace.SequenceCascade.ones_remainder_norm}\par{\sffamily\footnotesize Source SHA-256 prefix: 0a384fa460f07336.}}\par
\par\smallskip\noindent{\small\textbf{F} \path{SequenceCascade.lean}, line 270.\par\noindent\path{KomorebiSpace.SequenceCascade.geometric_reconstruction_error}\par{\sffamily\footnotesize Source SHA-256 prefix: 0a384fa460f07336.}}\par
\smallskip\noindent{\color{KomorebiMuted}\rule{\linewidth}{.3pt}}\par\end{minipage}\par\medskip
\par\noindent\begin{minipage}{\linewidth}\subsection*{KS30 / Gauge groupoid and section transport}
The action and groupoid instances are also read in this module. Composition order is retained.
\par\smallskip\noindent{\small\textbf{F} \path{GroupViewSystems.lean}, line 107.\par\noindent\path{KomorebiSpace.GroupViewSystem.twisted_coherent_iff}\par{\sffamily\footnotesize Source SHA-256 prefix: d2dd13e5f5cc2d78.}}\par
\par\smallskip\noindent{\small\textbf{F} \path{CocycleGaugeGroupoid.lean}, line 95.\par\noindent\path{KomorebiSpace.GroupViewSystem.gaugeSetoid}\par{\sffamily\footnotesize Source SHA-256 prefix: 93393322cd38df5e.}}\par
\smallskip\noindent{\color{KomorebiMuted}\rule{\linewidth}{.3pt}}\par\end{minipage}\par\medskip
\par\noindent\begin{minipage}{\linewidth}\subsection*{KS31 / Nonabelian extension obstruction}
A pointed based-gauge orbit set, not an assumed cohomology group or a general algorithm.
\par\smallskip\noindent{\small\textbf{F} \path{NonabelianExtension.lean}, line 96.\par\noindent\path{KomorebiSpace.GroupViewSystem.extensionOrbit_neutral_iff}\par{\sffamily\footnotesize Source SHA-256 prefix: 7834ff4193ad9a05.}}\par
\par\smallskip\noindent{\small\textbf{F} \path{NonabelianExtension.lean}, line 114.\par\noindent\path{KomorebiSpace.GroupViewSystem.extensionOrbit_independent}\par{\sffamily\footnotesize Source SHA-256 prefix: 7834ff4193ad9a05.}}\par
\smallskip\noindent{\color{KomorebiMuted}\rule{\linewidth}{.3pt}}\par\end{minipage}\par\medskip
\par\noindent\begin{minipage}{\linewidth}\subsection*{KS32 / Nonabelian extension torsor}
Parameters act by right multiplication on a chosen extension.
\par\smallskip\noindent{\small\textbf{F} \path{NonabelianExtensionFibers.lean}, line 51.\par\noindent\path{KomorebiSpace.GroupViewSystem.extensionParameters}\par{\sffamily\footnotesize Source SHA-256 prefix: 486c0bbabc34525c.}}\par
\par\smallskip\noindent{\small\textbf{F} \path{NonabelianExtensionFibers.lean}, line 107.\par\noindent\path{KomorebiSpace.GroupViewSystem.unique_extension_iff}\par{\sffamily\footnotesize Source SHA-256 prefix: 486c0bbabc34525c.}}\par
\smallskip\noindent{\color{KomorebiMuted}\rule{\linewidth}{.3pt}}\par\end{minipage}\par\medskip
\par\noindent\begin{minipage}{\linewidth}\subsection*{KS33 / S3 loop model}
Cocycle coherence is not the existence of a twisted global section.
\par\smallskip\noindent{\small\textbf{F} \path{CocycleExamples.lean}, line 205.\par\noindent\path{KomorebiSpace.CocycleExamples.s3_orbit_dichotomy}\par{\sffamily\footnotesize Source SHA-256 prefix: 6761e0443a1c1ad7.}}\par
\par\smallskip\noindent{\small\textbf{F} \path{CocycleExamples.lean}, line 169.\par\noindent\path{KomorebiSpace.CocycleExamples.actual_automorphisms_noncommute}\par{\sffamily\footnotesize Source SHA-256 prefix: 6761e0443a1c1ad7.}}\par
\par\smallskip\noindent{\small\textbf{F} \path{CocycleExamples.lean}, line 103.\par\noindent\path{KomorebiSpace.CocycleExamples.involution_trivial_iff}\par{\sffamily\footnotesize Source SHA-256 prefix: 6761e0443a1c1ad7.}}\par
\smallskip\noindent{\color{KomorebiMuted}\rule{\linewidth}{.3pt}}\par\end{minipage}\par\medskip
\par\noindent\begin{minipage}{\linewidth}\subsection*{KS34 / Effective groupoid descent}
Unframed gauge groupoids and composition-covered full patches. Uniqueness respects the specified identifications.
\par\smallskip\noindent{\small\textbf{F} \path{DescentEquivalence.lean}, line 104.\par\noindent\path{KomorebiSpace.GroupViewSystem.CocycleDescentData.descentEquivalence}\par{\sffamily\footnotesize Source SHA-256 prefix: c7ca73d617a578ed.}}\par
\par\smallskip\noindent{\small\textbf{F} \path{DescentComparison.lean}, line 79.\par\noindent\path{KomorebiSpace.GroupViewSystem.CocycleDescentData.Effectuation.exists_unique_comparison}\par{\sffamily\footnotesize Source SHA-256 prefix: afc5bf933284a648.}}\par
\smallskip\noindent{\color{KomorebiMuted}\rule{\linewidth}{.3pt}}\par\end{minipage}\par\medskip
\par\noindent\begin{minipage}{\linewidth}\subsection*{KS35 / Composition topology}
An actual Grothendieck topology, with the standard generated-sieve characterization.
\par\smallskip\noindent{\small\textbf{F} \path{CompositionSite.lean}, line 57.\par\noindent\path{KomorebiSpace.CompositionSite.topology}\par{\sffamily\footnotesize Source SHA-256 prefix: ac958295633fc50f.}}\par
\par\smallskip\noindent{\small\textbf{F} \path{CompositionSite.lean}, line 143.\par\noindent\path{KomorebiSpace.CompositionSite.ofArrows_covering_iff}\par{\sffamily\footnotesize Source SHA-256 prefix: ac958295633fc50f.}}\par
\smallskip\noindent{\color{KomorebiMuted}\rule{\linewidth}{.3pt}}\par\end{minipage}\par\medskip
\par\noindent\begin{minipage}{\linewidth}\subsection*{KS36 / Standard stackhood}
Target is mathlib Pseudofunctor.IsStack on the chosen composition site.
\par\smallskip\noindent{\small\textbf{F} \path{CompositionStack.lean}, line 178.\par\noindent\path{KomorebiSpace.GroupViewSystem.patch_isStack}\par{\sffamily\footnotesize Source SHA-256 prefix: 656bbbc1ddaba844.}}\par
\smallskip\noindent{\color{KomorebiMuted}\rule{\linewidth}{.3pt}}\par\end{minipage}\par\medskip
\par\noindent\begin{minipage}{\linewidth}\subsection*{KS37 / Equivalent view presentations}
Prescribed domains are isomorphism closed and coefficient systems are really pulled back.
\par\smallskip\noindent{\small\textbf{F} \path{SectionViewEquivalence.lean}, line 63.\par\noindent\path{KomorebiSpace.ViewSystem.viewEquivalence}\par{\sffamily\footnotesize Source SHA-256 prefix: 999bb5c576cefc7b.}}\par
\par\smallskip\noindent{\small\textbf{F} \path{RelativeSections.lean}, line 66.\par\noindent\path{KomorebiSpace.ViewSystem.relativeViewEquivalence}\par{\sffamily\footnotesize Source SHA-256 prefix: c28de5ab91b69828.}}\par
\par\smallskip\noindent{\small\textbf{F} \path{RelativeAmbiguity.lean}, line 43.\par\noindent\path{KomorebiSpace.GroupViewSystem.invisibleViewEquivalence}\par{\sffamily\footnotesize Source SHA-256 prefix: a24937a0816305f9.}}\par
\smallskip\noindent{\color{KomorebiMuted}\rule{\linewidth}{.3pt}}\par\end{minipage}\par\medskip
\par\noindent\begin{minipage}{\linewidth}\subsection*{KS38 / Identified-boundary realization groupoid}
Groupoid homotopy fiber construction. No new higher-category interface is asserted.
\par\smallskip\noindent{\small\textbf{F} \path{MarkedRealization.lean}, line 64.\par\noindent\path{KomorebiSpace.MarkedRealization.nonempty_iff}\par{\sffamily\footnotesize Source SHA-256 prefix: a3ad06525f28defd.}}\par
\par\smallskip\noindent{\small\textbf{F} \path{MarkedRealization.lean}, line 73.\par\noindent\path{KomorebiSpace.MarkedRealization.endomorphism_condition}\par{\sffamily\footnotesize Source SHA-256 prefix: a3ad06525f28defd.}}\par
\smallskip\noindent{\color{KomorebiMuted}\rule{\linewidth}{.3pt}}\par\end{minipage}\par\medskip
\par\noindent\begin{minipage}{\linewidth}\subsection*{KS39 / Marked transport across a square}
The square commutes through a specified natural isomorphism and the horizontal functors are equivalences.
\par\smallskip\noindent{\small\textbf{F} \path{MarkedTransport.lean}, line 117.\par\noindent\path{KomorebiSpace.MarkedRealization.transportEquivalence}\par{\sffamily\footnotesize Source SHA-256 prefix: 99fc0cf03cfdbdf9.}}\par
\par\smallskip\noindent{\small\textbf{F} \path{MarkedTransport.lean}, line 128.\par\noindent\path{KomorebiSpace.MarkedRealization.transportHomEquiv}\par{\sffamily\footnotesize Source SHA-256 prefix: 99fc0cf03cfdbdf9.}}\par
\smallskip\noindent{\color{KomorebiMuted}\rule{\linewidth}{.3pt}}\par\end{minipage}\par\medskip
\par\noindent\begin{minipage}{\linewidth}\subsection*{KS40 / Transported path telescope}
Known bounded invertible transports. Path gain is not replaced by length without contractivity.
\par\smallskip\noindent{\small\textbf{F} \path{AmplifiedTransport.lean}, line 151.\par\noindent\path{KomorebiSpace.AmplifiedTransport.Connection.walkSum_defect}\par{\sffamily\footnotesize Source SHA-256 prefix: da79596a2f437a07.}}\par
\par\smallskip\noindent{\small\textbf{F} \path{AmplifiedTransport.lean}, line 239.\par\noindent\path{KomorebiSpace.AmplifiedTransport.Connection.walk_error}\par{\sffamily\footnotesize Source SHA-256 prefix: da79596a2f437a07.}}\par
\smallskip\noindent{\color{KomorebiMuted}\rule{\linewidth}{.3pt}}\par\end{minipage}\par\medskip
\par\noindent\begin{minipage}{\linewidth}\subsection*{KS41 / Holonomy classification}
A bounded linear equivalence. It is an isometry only in the isometric subclass.
\par\smallskip\noindent{\small\textbf{F} \path{AmplifiedHolonomy.lean}, line 174.\par\noindent\path{KomorebiSpace.AmplifiedTransport.Rooted.parameterEquiv}\par{\sffamily\footnotesize Source SHA-256 prefix: f46a0ba9482bc51c.}}\par
\par\smallskip\noindent{\small\textbf{F} \path{AmplifiedHolonomy.lean}, line 104.\par\noindent\path{KomorebiSpace.AmplifiedTransport.Rooted.fixedSpace_route_independent}\par{\sffamily\footnotesize Source SHA-256 prefix: f46a0ba9482bc51c.}}\par
\smallskip\noindent{\color{KomorebiMuted}\rule{\linewidth}{.3pt}}\par\end{minipage}\par\medskip
\par\noindent\begin{minipage}{\linewidth}\subsection*{KS42 / Orthogonal exact repair}
Requires the declared Hilbert geometry, not closed residual range.
\par\smallskip\noindent{\small\textbf{M} \path{HilbertKernelProjection.lean}, line 245.\par\noindent\path{KomorebiSpace.HilbertKernelProjection.unconditional_exact_repair}\par{\sffamily\footnotesize Source SHA-256 prefix: 5a30a4de9c18ed34.}}\par
\par\smallskip\noindent{\small\textbf{M} \path{HilbertKernelProjection.lean}, line 120.\par\noindent\path{KomorebiSpace.HilbertKernelProjection.projection_best}\par{\sffamily\footnotesize Source SHA-256 prefix: 5a30a4de9c18ed34.}}\par
\par\smallskip\noindent{\small\textbf{M} \path{HilbertKernelProjection.lean}, line 134.\par\noindent\path{KomorebiSpace.HilbertKernelProjection.projection_unique}\par{\sffamily\footnotesize Source SHA-256 prefix: 5a30a4de9c18ed34.}}\par
\smallskip\noindent{\color{KomorebiMuted}\rule{\linewidth}{.3pt}}\par\end{minipage}\par\medskip
\par\noindent\begin{minipage}{\linewidth}\subsection*{KS43 / Closed range equivalence}
Banach target in the stated theorem. The explicit infinite diagonal illustration is written, not newly compiled.
\par\smallskip\noindent{\small\textbf{M} \path{HilbertKernelProjection.lean}, line 210.\par\noindent\path{KomorebiSpace.HilbertKernelProjection.closed_range_iff_bound}\par{\sffamily\footnotesize Source SHA-256 prefix: 5a30a4de9c18ed34.}}\par
\par\smallskip\noindent{\small\textbf{M} \path{HilbertKernelProjection.lean}, line 219.\par\noindent\path{KomorebiSpace.HilbertKernelProjection.closed_range_iff_near_exact}\par{\sffamily\footnotesize Source SHA-256 prefix: 5a30a4de9c18ed34.}}\par
\par\smallskip\noindent{\small\textbf{M} \path{HilbertKernelProjection.lean}, line 163.\par\noindent\path{KomorebiSpace.HilbertKernelProjection.gain_characterization}\par{\sffamily\footnotesize Source SHA-256 prefix: 5a30a4de9c18ed34.}}\par
\smallskip\noindent{\color{KomorebiMuted}\rule{\linewidth}{.3pt}}\par\end{minipage}\par\medskip
\par\noindent\begin{minipage}{\linewidth}\subsection*{KS44 / Amplified holonomy gain}
The finite-dimensional result is in the foundation archive. The infinite-dimensional necessity is in a separately recovered source, not claimed under the focused archive replay.
\par\smallskip\noindent{\small\textbf{F} \path{AmplifiedRepair.lean}, line 204.\par\noindent\path{KomorebiSpace.AmplifiedTransport.Rooted.orthogonalRepair_error}\par{\sffamily\footnotesize Source SHA-256 prefix: e4009d19b75ab9cf.}}\par
\par\smallskip\noindent{\small\textbf{S} \path{HilbertHolonomyRepair.lean}, line 194.\par\noindent\path{closed_range_iff_near_parallel}\par{\sffamily\footnotesize Source SHA-256 prefix: 5baf9b3a40b4fdd8.}}\par
\smallskip\noindent{\color{KomorebiMuted}\rule{\linewidth}{.3pt}}\par\end{minipage}\par\medskip
\par\noindent\begin{minipage}{\linewidth}\subsection*{KS45 / Flat transport amplification}
Both directed edge orientations use the original maximum norm.
\par\smallskip\noindent{\small\textbf{F} \path{AmplifiedExamples.lean}, line 139.\par\noindent\path{KomorebiSpace.AmplifiedTransport.Examples.exponential_amplification}\par{\sffamily\footnotesize Source SHA-256 prefix: e95779e1b481e79e.}}\par
\par\smallskip\noindent{\small\textbf{F} \path{AmplifiedExamples.lean}, line 148.\par\noindent\path{KomorebiSpace.AmplifiedTransport.Examples.length_only_bound_false}\par{\sffamily\footnotesize Source SHA-256 prefix: e95779e1b481e79e.}}\par
\smallskip\noindent{\color{KomorebiMuted}\rule{\linewidth}{.3pt}}\par\end{minipage}\par\medskip
\par\noindent\begin{minipage}{\linewidth}\subsection*{KS46 / Two-stage joint projection}
Second projection is inside the locally exact subspace.
\par\smallskip\noindent{\small\textbf{M} \path{QuantitativeDescent.lean}, line 61.\par\noindent\path{KomorebiSpace.QuantitativeDescent.repair_eq_projection}\par{\sffamily\footnotesize Source SHA-256 prefix: 0a35307bbdb5fdc3.}}\par
\par\smallskip\noindent{\small\textbf{M} \path{QuantitativeDescent.lean}, line 121.\par\noindent\path{KomorebiSpace.QuantitativeDescent.error_squared}\par{\sffamily\footnotesize Source SHA-256 prefix: 0a35307bbdb5fdc3.}}\par
\smallskip\noindent{\color{KomorebiMuted}\rule{\linewidth}{.3pt}}\par\end{minipage}\par\medskip
\par\noindent\begin{minipage}{\linewidth}\subsection*{KS47 / Joint closed range}
Closed local range is retained. The printed bound is a sufficient, not necessarily optimal, gain.
\par\smallskip\noindent{\small\textbf{M} \path{DescentClosedRange.lean}, line 93.\par\noindent\path{KomorebiSpace.QuantitativeDescent.joint_closed_iff_matching_closed}\par{\sffamily\footnotesize Source SHA-256 prefix: 394f9a40603290cd.}}\par
\par\smallskip\noindent{\small\textbf{M} \path{DescentClosedRange.lean}, line 70.\par\noindent\path{KomorebiSpace.QuantitativeDescent.stable_repair}\par{\sffamily\footnotesize Source SHA-256 prefix: 394f9a40603290cd.}}\par
\smallskip\noindent{\color{KomorebiMuted}\rule{\linewidth}{.3pt}}\par\end{minipage}\par\medskip
\par\noindent\begin{minipage}{\linewidth}\subsection*{KS48 / Two-sided projector motion}
Both independent residual certificates are supplied. No inclusion of kernels is assumed.
\par\smallskip\noindent{\small\textbf{M} \path{MovingKernelBounds.lean}, line 95.\par\noindent\path{KomorebiSpace.MovingKernel.projection_motion}\par{\sffamily\footnotesize Source SHA-256 prefix: 0edf4cbd9a1b1063.}}\par
\par\smallskip\noindent{\small\textbf{M} \path{MovingKernelBounds.lean}, line 164.\par\noindent\path{KomorebiSpace.MovingKernel.projection_lipschitz}\par{\sffamily\footnotesize Source SHA-256 prefix: 0edf4cbd9a1b1063.}}\par
\smallskip\noindent{\color{KomorebiMuted}\rule{\linewidth}{.3pt}}\par\end{minipage}\par\medskip
\par\noindent\begin{minipage}{\linewidth}\subsection*{KS49 / Derived perturbed certificate}
Banach target, supplied split surjectivity, and strict a times epsilon less than one.
\par\smallskip\noindent{\small\textbf{M} \path{SplitKernelPerturbation.lean}, line 92.\par\noindent\path{KomorebiSpace.SplitKernel.Certificate.perturbed_right_inverse}\par{\sffamily\footnotesize Source SHA-256 prefix: 1e60514ffacf125c.}}\par
\par\smallskip\noindent{\small\textbf{M} \path{SplitKernelPerturbation.lean}, line 98.\par\noindent\path{KomorebiSpace.SplitKernel.Certificate.perturbed_section_bound}\par{\sffamily\footnotesize Source SHA-256 prefix: 1e60514ffacf125c.}}\par
\par\smallskip\noindent{\small\textbf{M} \path{SplitKernelPerturbation.lean}, line 140.\par\noindent\path{KomorebiSpace.SplitKernel.Certificate.projection_motion}\par{\sffamily\footnotesize Source SHA-256 prefix: 1e60514ffacf125c.}}\par
\smallskip\noindent{\color{KomorebiMuted}\rule{\linewidth}{.3pt}}\par\end{minipage}\par\medskip
\par\noindent\begin{minipage}{\linewidth}\subsection*{KS50 / Transport between moving kernels}
Transport depends on the splitting and is not generally orthogonal projection.
\par\smallskip\noindent{\small\textbf{M} \path{SplitKernelPerturbation.lean}, line 216.\par\noindent\path{KomorebiSpace.SplitKernel.Certificate.kernelEquiv}\par{\sffamily\footnotesize Source SHA-256 prefix: 1e60514ffacf125c.}}\par
\par\smallskip\noindent{\small\textbf{M} \path{SplitKernelPerturbation.lean}, line 226.\par\noindent\path{KomorebiSpace.SplitKernel.Certificate.kernel_transport_error}\par{\sffamily\footnotesize Source SHA-256 prefix: 1e60514ffacf125c.}}\par
\smallskip\noindent{\color{KomorebiMuted}\rule{\linewidth}{.3pt}}\par\end{minipage}\par\medskip
\par\noindent\begin{minipage}{\linewidth}\subsection*{KS51 / Exact moving-line example}
The all-real projector formula is separate from the strict-radius certificate interface.
\par\smallskip\noindent{\small\textbf{M} \path{MovingKernelExamples.lean}, line 91.\par\noindent\path{KomorebiSpace.MovingKernelExamples.projection_formula}\par{\sffamily\footnotesize Source SHA-256 prefix: 8ce528290f97bfde.}}\par
\par\smallskip\noindent{\small\textbf{M} \path{MovingKernelExamples.lean}, line 142.\par\noindent\path{KomorebiSpace.MovingKernelExamples.motion_opNorm_exact}\par{\sffamily\footnotesize Source SHA-256 prefix: 8ce528290f97bfde.}}\par
\par\smallskip\noindent{\small\textbf{M} \path{MovingKernelExamples.lean}, line 199.\par\noindent\path{KomorebiSpace.MovingKernelExamples.kernel_transport_formula}\par{\sffamily\footnotesize Source SHA-256 prefix: 8ce528290f97bfde.}}\par
\smallskip\noindent{\color{KomorebiMuted}\rule{\linewidth}{.3pt}}\par\end{minipage}\par\medskip
\par\noindent\begin{minipage}{\linewidth}\subsection*{KS52 / Rank loss and radius endpoint}
The rank-changing diagonal pair is a written elementary example. Source correspondence covers the general loss estimate and the radius boundary.
\par\smallskip\noindent{\small\textbf{M} \path{MovingKernelBounds.lean}, line 148.\par\noindent\path{KomorebiSpace.MovingKernel.lost_exact_unit_forces_gain}\par{\sffamily\footnotesize Source SHA-256 prefix: 0edf4cbd9a1b1063.}}\par
\par\smallskip\noindent{\small\textbf{M} \path{MovingKernelExamples.lean}, line 227.\par\noindent\path{KomorebiSpace.MovingKernelExamples.closed_radius_failure}\par{\sffamily\footnotesize Source SHA-256 prefix: 8ce528290f97bfde.}}\par
\smallskip\noindent{\color{KomorebiMuted}\rule{\linewidth}{.3pt}}\par\end{minipage}\par\medskip
\par\noindent\begin{minipage}{\linewidth}\subsection*{KS53 / Motion of quantitative descent}
Uses the original two-stage repair through its proved identification with the joint projector.
\par\smallskip\noindent{\small\textbf{M} \path{MovingDescent.lean}, line 52.\par\noindent\path{KomorebiSpace.MovingDescent.repair_motion}\par{\sffamily\footnotesize Source SHA-256 prefix: 6820862031ad7094.}}\par
\smallskip\noindent{\color{KomorebiMuted}\rule{\linewidth}{.3pt}}\par\end{minipage}\par\medskip
\par\noindent\begin{minipage}{\linewidth}\subsection*{KS54 / Finite prefixes without a point}
Additional written counterexample with a complete elementary proof. No new Lean declaration or universal computational proof is claimed.
\par\smallskip\noindent\textbf{W.} The proof is written in the monograph. No new Lean declaration or compiler verification is claimed.
\smallskip\noindent{\color{KomorebiMuted}\rule{\linewidth}{.3pt}}\par\end{minipage}\par\medskip


\clearpage\section{Reproduction without changing the evidence}
\subsection{Current publication checks}
From the root of the extracted publication package, run:
\begin{quote}\small\ttfamily
python3 checks/verify\_sources.py\par
python3 checks/build\_concordance.py\par
python3 checks/exact\_examples.py\par
python3 checks/build\_lean\_guide.py
\end{quote}
The supplied build script runs these steps, invokes the recovered focused exact checker, and compiles the three PDFs. It uses a TeX Live installation with the packages declared in the shared preamble. It is not a Lean installation script.

Keep current outputs separate from inherited reports. A new run may change timestamps in a current report. It must not modify the archived reports to make them appear newly executed.

\subsection{Replaying the original Lean projects}
Extract each original source ZIP into its own directory. Supply the pinned portable Lean runtime and its matching dependency cache. The original README in each project gives the precise verifier interface. The foundation project supports a full verification run and a separate construction audit. The focused project requires a fresh build name and compiles its entire supplied closure.

The runtime is not embedded in this publication. The source tree and the pinned dependency records are included, but restoring an external compiler and cache is a separate operation. Neither a missing runtime nor a process timeout is a proof failure. Conversely, a process that did not complete must not be marked PASS.

\subsection{A future consolidated replay}
A combined Lean release should start by selecting the canonical version of every shared module and comparing any differences. It should preserve the old source archives, make every import explicit, and map the printed hypotheses to the chosen declarations. Only then should it execute a fresh full build, axiom audit, and negative controls on the integrated tree.

The present publication does not claim that this integration has already occurred. Its two complete recovered source trees are instead reproducible, separately identified inputs to that next task.

\section{Editorial review decisions}
Two provisional declaration names did not match the source inventory during construction of the concordance. They were replaced by the actual names \code{onLimit_comp} and \code{closed_radius_failure}; the resolver then passed all entries. A malformed norm token in the draft and a long hash line that exceeded the text width were also corrected during the PDF review.

The monograph retains the order of nonabelian multiplication, the difference between a prescribed boundary and a literal representative, the closed-range hypotheses of stable repair, and the distinction between an anchor projection and a global nearest-point problem. It also states that the strict set of coherent families can be empty even when every view is nonempty.

\end{document}
